\section{Experimental Setup}
\label{sec:setup}

\subsection{Capability-Conditioned Loss}
\label{sec:loss}

We measure three capabilities, mathematical reasoning, code
generation, and multi-hop question answering, each through a fixed
probe set $\probeset{c}$ of prompt and reference pairs drawn once
from held-out splits of MATH-500
\citep{hendrycks2021math,lightman2023verify}, MBPP
\citep{austin2021program}, and 2WikiMultihopQA
\citep{ho2020constructing}. Each probe set fixes one evaluation distribution $j$, so the measured
quantity is $L_{c,j}$; we write $L_c$ with the fixed distribution
implied and test in Appendix~\ref{sec:app_measure} how far a response
carries to other distributions of the same capability. The observable is the
token-weighted conditional cross-entropy of the reference completion,
\begin{equation}
\Lc{c}(w)\;=\;\frac{1}{T_c}\sum_{(x,y)\in\probeset{c}}\;\sum_{t=1}^{|y|}
-\log p_w\!\left(y_t\mid x,y_{<t}\right),\qquad
T_c=\sum_{(x,y)\in\probeset{c}}|y|,
\label{eq:taskloss}
\end{equation}
with prompt tokens masked, 1024-token truncation, the model's own
tokenizer, and the same 64 measurement probes per capability for
every state compared within a model. We define the reference model $M_0$ as the checkpoint before the
compression intervention; the definition accommodates dense and
sparsely activated architectures, while our controlled experiments
focus on dense models. Responses are within-model differences on
identical token sequences in native nats per token: for pruning and
quantization the damage
$\Delta L_c=L_c(\text{compressed})-L_c(M_0)$, for distillation the
signed transfer $\delta_c=L_c(S_{\mathrm{KD}})-L_c(S_0)$ of a student
relative to its own initial state $S_0$, which is the reference
model of that arm; the teacher never plays the role of the source. The observable is the likelihood of a fixed reference on a fixed
distribution: a QA loss can fall while exact-match accuracy does not
rise, and a gain on one QA distribution need not appear on another;
absolute levels of different tokenizers are never placed on one scale
(Appendix~\ref{sec:app_measure}).

\subsection{Models and Compression Settings}
\label{sec:models}

A \emph{heterogeneous panel} of twelve finished public models from
five series (Qwen3, Gemma-3, Gemma-4, Muse, OLMo-3; 0.6B to 32B)
establishes the phenomena and the limits of cross-family transfer
(Appendix~\ref{sec:app_hetero}). A \emph{controlled panel} of Pythia checkpoints
\citep{biderman2023pythia} (160M, 410M, 1.4B at steps 16k, 64k, 143k) varies the
initial state along parameter count and pretraining stage with
architecture and data fixed; it is the development set for pruning and quantization and the controlled transfer panel, while the distillation relations are developed on Gemma-3 student trajectories (\S\ref{sec:law_distill}). \emph{New states and configurations} that entered no fit (Pythia
160M to 6.9B at further steps, new densities, bit-widths, and group
sizes, Gemma-3 students at new pools; Appendix~\ref{sec:app_details})
provide the validation.

The three operations are held fixed by protocol: global magnitude
pruning of all transformer weight matrices to a retained density $d$;
symmetric round-to-nearest weight quantization of the same matrices,
per output channel at bit-width $b$ or with contiguous input groups of
$g$ weights sharing one scale;
and supervised fine-tuning of an existing student on black-box
teacher traces with strict LoRA adapters \citep{hu2022lora} from a
pool of $U$ traces per domain, processed for a fixed budget. The inputs admitted to every predictor are the parameter count $N_0$
of the transformer weight matrices of $M_0$ (total and per-token
active counts coincide for the dense models studied here), the
pretraining token count $D_0$ where disclosed, and the reference
capability losses $L_{0,c}=L_c(M_0)$ measured once (the initial
student for distillation). Every prediction is written as
\begin{equation}
\widehat L_{m,c}\;=\;L^{\mathrm{obs}}_{0,c}\;+\;
\widehat\Delta_{m,c}(\mathbf x,\boldsymbol\theta_m),\qquad
\mathbf x=(N_0,D_0,L_{0,c}),
\label{eq:interface}
\end{equation}
with $\boldsymbol\theta_m$ the setting of method $m$ (density,
bit-width and group size, or pool and budget). The function $\widehat\Delta_{m,c}$ is method-specific; the arms
share the endpoint and the evaluation protocol.

\subsection{Fitting and Evaluation}
\label{sec:eval}

Each relation is fit on development states and settings only and is
judged on unseen settings (outside the fitted set) and unseen states
(a new size, stage, or student). On the development panel we use leave-one-stage-out and
leave-one-source-out cross-validation; for new states every
candidate's forms, coefficients, and predictions are committed before
the outcome is measured, and a prediction keeps its provenance label
whatever the outcome. The reference anchor $L_{0,c}$ is the only measured input. Error is the mean
absolute error of $\widehat\Delta$ in nats per token per capability;
every candidate is compared on the same cells with source-free
references (zero change, a per-capability constant, a strength-only
curve, or the median development curve), with same-input alternatives,
and with the no-$D_0$ ablation, and the strongest baseline per
capability is reported and marked as chosen after the fact. Freeze
commits, provenance codes, fold rules, and bootstrap units are in
Appendix~\ref{sec:app_registration}.
